\documentclass[11pt,a4paper]{article}
\usepackage[margin=25mm]{geometry}
\usepackage{amsmath,amssymb,amsthm,hyperref,array}
\setlength{\emergencystretch}{2em}
\newtheorem{proposition}{Proposition}
\newcommand{\gen}[1]{\langle#1\rangle}
\newcommand{\Alt}{\operatorname{Alt}}
\title{Completing the dihedral attachment branch}
\author{Human Compression project: bounded L044 proof component}
\date{}
\begin{document}\maketitle
Let $G=\gen{\rho_0,\ldots,\rho_{r-1}}\leq A_n$, $n\geq12$,
be a transitive string C-group whose maximal omissions are
intransitive. Suppose a unique eligible $2$-edge joins the
$G_2$-orbits $A=\{1,\ldots,5\}$ and $X$, with
\[
 \rho_0=l=(1\ 2)(3\ 4),\qquad
 \rho_1=s=(2\ 3)(4\ 5).
\]
Assume $K=G_{>2}$ is supported on and transitive on $X$;
write $q=|X|$ and $d=r-3$. The crossing endpoint in $A$ is
necessarily $5$, the sole fixed point of $l$.
We assume the alternating rank bound at every strictly smaller
degree at least $12$. The accepted proper-primitive,
anchored-imprimitive and small anchored-interval results are
paid supporting inputs, not a blanket unanchored rank bound.

\begin{proposition}
Either $r\leq\lfloor(n-1)/2\rfloor$, or the actual prefix
$\gen{\rho_0,\rho_1,\rho_2}$, with its ordered generators,
is exactly source Table~2 row~(2) or row~(3) on six points
and fixes the remaining points. The dual statement holds.
\end{proposition}
\begin{proof}
The tail $K$ has a genuine even string C-marking. Every
maximal omission in $K$ is intransitive on $X$: otherwise
adjoining the head, transitive on $A$, and the crossing
involution would make the corresponding global omission
transitive. All its generators except the first fix the
crossing endpoint in $X$, by far commutation. Its diagram is
connected, since a commuting marked component away from
that endpoint would fix it and hence, by transitivity and
commutation, fix all of $X$.

If $q<12$ and $d\geq(q-1)/2$, the accepted small primitive
and anchored-imprimitive interval results leave only source
Table~2 rows~(1)--(3). Their support degrees are five or six,
which would give $n\leq11$. Thus $d<(q-1)/2$ for $q<12$.
For $q\geq12$, the accepted proper-primitive and anchored-
imprimitive bounds give $d\leq\lfloor(q-2)/2\rfloor$
unless $K=\Alt(X)$. In all these cases
\[
       r=d+3\leq\left\lfloor\frac{q+4}{2}\right\rfloor
                  =\left\lfloor\frac{n-1}{2}\right\rfloor.
\]
In the remaining case the explicit smaller-degree induction
hypothesis gives $d\leq\lfloor(q-1)/2\rfloor$.
A failure of the asserted bound therefore requires
\[
 q\geq13\text{ odd},\qquad K=\Alt(X),\qquad
 d=(q-1)/2.
\]

Name the crossing endpoint in $X$ as $6$. The unique cross
transposition and $[l,\rho_2]=1$ give
\[
             \rho_2=(5\ 6)u v,
\]
where $u$ is an involution in $C_{S_{\{1,2,3,4\}}}(l)$,
and $v$ is an involution on $X\setminus\{6\}$.
Their signs are opposite. The centralizer preserves the two
pairs of $l$ and has order eight; its involutions, including
the identity, are exactly
\[
 1,\ (1\ 2),\ (3\ 4),\ (1\ 2)(3\ 4),\
 (1\ 3)(2\ 4),\ (1\ 4)(2\ 3).
\]
We dispatch all six literal choices, without replacing a
marked tuple by an abstract group of the same order.

If $u=1$, then $R=\gen{\rho_2,K}$ is an actual even
subgroup supported on $X\cup\{5\}$. It is the full
$\Alt(X\cup\{5\})$, by overlapping natural alternating
supports. Its full IP is inherited from $G$, since no
projection was taken. Its rank is $d+1=(q+1)/2$ and its
degree $q+1$ is strictly smaller than $n=q+5$ and at least
$14$. This contradicts the stated induction hypothesis.

If $u=(1\ 2)(3\ 4)$ or $u=(1\ 3)(2\ 4)$, the edges of
$s$ and $u$ connect all five points of $A$. Together with
$K$ on $X$ and the crossing edge, the actual subgroup
$G_0=\gen{s,\rho_2,K}$ is transitive on $\Omega$.
This is forbidden. For the first choice the connecting path
is $1,2,3,4,5$; for the second it is $1,3,2,4,5$.

If $u=(1\ 4)(2\ 3)$, conjugate by the permutation
$(1\ 4\ 2\ 3)$, fixing $5$ and $X$. The first three
marks become
\[
 (1\ 2)(3\ 4),\quad(2\ 5)(1\ 3),\quad
 (1\ 3)(2\ 4)(5\ 6)v.
\]
The complete dihedral--alternating obstruction applies and
excludes this case. In particular the decorated six-point
$S_5$ possibility is not silently identified with either
of the source's $A_5$ rows.

It remains to take $u=(1\ 2)$ or $u=(3\ 4)$.
Projection of $P=\gen{l,s,\rho_2}$ to $\{1,\ldots,6\}$
has the literal generators of accepted source Table~2
row~(2) or row~(3), respectively. Both are the perfect
group $A_5$ in that precise degree-six action. The exterior
projection of $P$ lies in $\gen v$. If $v\ne1$ and the
six-point projection were injective, the exterior action
would define a nontrivial homomorphism $A_5\to C_2$.
Consequently $v\in P$. But $v$ is even, so belongs to
$K=\Alt(X)$, while full ambient IP gives $P\cap K=1$.
Thus $v=1$. The two remaining actual ordered prefixes are
exactly the two stated source rows. Reversing the full
string gives the dual conclusion.
\end{proof}

\paragraph{Source, dependencies and correction scope.}
This is the complete dihedral attachment argument for
Cameron--Fernandes--Leemans--Mixer, \emph{Highest rank of
a polytope for $A_n$}, arXiv:1605.09173v1,
Proposition~5.14, physical pages~28--29, in the supported
interval setting of its Section~5.2. It reuses the source's
rank-deficit and sesqui-extension strategy. The six-choice
normalization is stated explicitly and the previously
uncovered $S_5$ decoration is excluded by the new full-IP
obstruction. The original conclusion is retained; no
counterexample to that conclusion is asserted.

Both source rows and their exact $A_5$ identifications are
already covered by the accepted Table~2 packet; its complete
required census and source transcription remain paid
inputs. No census is repeated here. The alternating rank
step is an explicit strong-induction hypothesis, not a claim
that all remaining L044 branches have already been proved.
The exceptional prefixes, their later source cases, the
connected-square branch and the full source assembly remain
separate. No new formal, original or question credit follows
from this component.
\end{document}
